
\documentclass[12pt,a4paper]{article}
\usepackage[utf8]{inputenc}
\usepackage[T1]{fontenc}
\usepackage[a4paper,margin=2.2cm]{geometry}
\usepackage{setspace}
\usepackage{xcolor}
\usepackage{booktabs}
\usepackage{longtable}
\usepackage{array}
\usepackage{listings}
\usepackage{hyperref}
\usepackage{enumitem}
\usepackage{caption}
\setstretch{1.12}
\setlength{\parindent}{0pt}
\setlength{\parskip}{6pt}
\setlist[itemize]{topsep=3pt,itemsep=2pt,parsep=0pt}
\setlist[enumerate]{topsep=3pt,itemsep=2pt,parsep=0pt}
\hypersetup{colorlinks=true,linkcolor=blue!55!black,citecolor=blue!55!black,urlcolor=blue!55!black}
\definecolor{softgray}{RGB}{245,245,245}
\definecolor{linegray}{RGB}{70,70,70}
\lstdefinestyle{jsonstyle}{basicstyle=\ttfamily\small,backgroundcolor=\color{softgray},frame=single,rulecolor=\color{linegray},breaklines=true,showstringspaces=false,tabsize=2}
\lstset{style=jsonstyle}
\newcommand{\code}[1]{\texttt{#1}}
\newcommand{\formalcover}[4]{%
\begin{titlepage}
\centering
\vspace*{1.2cm}
{\Large Universitat Autonoma de Barcelona}\\[4pt]
{\large Master Degree in Modelling for Science and Engineering}\\[4pt]
{\large Institution / Lab: IIIA-CSIC}\\[2.0cm]
{\Huge\bfseries #1}\\[0.35cm]
{\Large #2}\\[1.5cm]
{\large Formal architecture and contract section for technical review.}\\[2.0cm]
\begin{tabular}{>{\bfseries}p{0.28\linewidth}p{0.58\linewidth}}
Document type: & Architecture and contract walkthrough\\
Academic year: & 2025--2026\\
Date: & #3\\
Scope: & #4\\
Repository: & \url{https://github.com/ieverythng/nao-ros4hri-bridge}\\
\end{tabular}
\vfill
\end{titlepage}}

\begin{document}
\formalcover{Technical Memory Draft Baseline}{End-to-End Planner Dialogue Flow}{2026-06-01}{Runtime flow, JSON contracts, replanning lineage, and dialogue relay behavior.}

\section{End-to-End User to Planner Dialogue Flow}

\subsection{Purpose}

This document is the supervisor-facing and thesis-facing walkthrough for the active planner stack. It is payload-first, but it now separates two world-state views:

\begin{itemize}
\item \code{grounded\_context}: compact scene graph for chatbot/planner LLM calls.
\item execution evidence: live skill-local data from \code{/scene/summary}, KB triples, tracked people, action results, and detector metadata.
\end{itemize}

The design follows the current neurosymbolic planning pattern used by SayPlan, ProgPrompt, RePLan, and scene-graph replanning work: pass only the task-relevant symbolic subgraph to the LLM, keep symbolic validation outside the LLM, and let execution skills refresh live evidence before claiming results.

Companion references:

\begin{itemize}
\item \code{docs/contracts.md}
\item \code{docs/current\_workflow.md}
\item \code{docs/plans/planner\_grounding\_moe\_contract\_2026-06-01.md}
\item \code{docs/plans/planner\_replan\_lineage\_adr.md}
\end{itemize}



\subsection{End-to-End Runtime Flow}

\begin{lstlisting}[caption={Diagram source from canonical Markdown}]
sequenceDiagram
    participant U as User
    participant C as chatbot_llm
    participant O as nao_orchestrator
    participant P as planner_llm
    participant S as AB=1 Skills
    participant K as KB + Scene Sources
    participant D as dialogue_manager

    U->>C: Natural language request
    C->>K: Query KB rows + scene summary
    C->>C: Project compact grounded_context
    C->>O: /nao_orchestrator/planner_request
    O->>O: PlannerGate admission
    O->>P: /planner/request
    P->>P: Build prompt + validate model plan
    P->>O: /intents
    O->>S: Dispatch executable steps
    S->>K: Requery live evidence
    S-->>O: Typed skill result
    O->>P: /planner/execution_feedback
    P->>O: /planner/dialogue_act when needed
    O->>D: /nao_orchestrator/planner_dialogue_act
    D-->>U: Spoken response
\end{lstlisting}

Ownership rule:

\begin{itemize}
\item \code{chatbot\_llm} owns user-facing interpretation and compact grounding projection.
\item \code{planner\_llm} owns planning, validation, retry, replan, and dialogue-act intent.
\item \code{nao\_orchestrator} owns deterministic admission, dispatch, and feedback publication.
\item AB=1 skills own live world interaction and evidence collection.
\item \code{dialogue\_manager} owns speech realization.
\end{itemize}



\subsection{Contract Map}

\begin{longtable}{p{0.23\linewidth}p{0.23\linewidth}p{0.23\linewidth}p{0.23\linewidth}}
Contract & Channel & Producer & Consumer\\
\midrule
Chatbot turn output & internal turn result JSON & \code{chatbot\_llm} & \code{chatbot\_llm} handoff path\\
Compact grounding & \code{grounded\_context} & \code{chatbot\_llm} & \code{chatbot\_llm}, \code{planner\_llm}\\
Planner request envelope & \code{/nao\_orchestrator/planner\_request}, \code{/planner/request} & \code{chatbot\_llm}, admitted by \code{nao\_orchestrator} & \code{planner\_llm}\\
Planner output & \code{/intents} & \code{planner\_llm} & \code{nao\_orchestrator}\\
Execution feedback & \code{/planner/execution\_feedback} & \code{nao\_orchestrator} & \code{planner\_llm}\\
Planner dialogue act & \code{/planner/dialogue\_act} & \code{planner\_llm} & \code{nao\_orchestrator} relay to \code{dialogue\_manager}\\
Raw scene/KB evidence & \code{/scene/summary}, \code{/kb/query}, skill action results & scene/KB/skills & skill-local execution paths\\
\end{longtable}

Removed request seams:

\begin{itemize}
\item \code{requested\_plan} is not part of the active planner request contract.
\item \code{interaction\_mode} is not part of the active planner request contract.
\item Legacy fields such as \code{goal\_token}, \code{world\_model\_snapshot}, and \code{world\_model\_text} remain removed.
\end{itemize}



\subsection{Chatbot Turn Output}

\begin{lstlisting}
{
  "verbal_ack": "Okay, I will do that.",
  "route": "execution",
  "confidence": 0.82,
  "user_intent": {
    "type": "look_at",
    "goal": "look at the person",
    "goal_text": "look at the person",
    "scene_targets": ["person"],
    "request_kind": "new_goal"
  }
}
\end{lstlisting}

Variables:

\begin{itemize}
\item \code{verbal\_ack}: optional immediate acknowledgement text.
\item \code{route}: \code{dialogue | knowledge\_query | execution}.
\item \code{confidence}: route/intent confidence in \code{[0.0, 1.0]}.
\item \code{user\_intent.type}: normalized intent label.
\item \code{user\_intent.goal} / \code{goal\_text}: planner-facing objective text.
\item \code{user\_intent.scene\_targets}: compact target labels or entity IDs.
\item \code{user\_intent.request\_kind}: request transition class.
\end{itemize}



\subsection{Compact Grounded Context}

The LLM-facing world state is a concise scene graph, not a raw RDF dump.

\begin{lstlisting}
{
  "grounded_context": {
    "entities": [
      {
        "id": "cup_jrjic",
        "label": "cup",
        "kind": "object",
        "class": "Cup",
        "visible": true,
        "relations": [
          {"predicate": "dbp:color", "object": "blue"},
          {"predicate": "oro:isOn", "object": "table_1"}
        ]
      },
      {
        "id": "anonymous_person_ehfbf",
        "label": null,
        "kind": "person",
        "class": "Human",
        "visible": true,
        "relations": []
      }
    ]
  }
}
\end{lstlisting}

Policy:

\begin{itemize}
\item \code{id} is the stable entity handle.
\item \code{label} is the human-readable class/name hint and may be \code{null} for anonymous people.
\item \code{kind} is \code{object} or \code{person}.
\item \code{class} is the compact semantic class.
\item \code{visible} is the current visual grounding flag.
\item \code{relations} only carries prioritized semantic predicates for LLM reasoning when they add information beyond \code{class}.
\end{itemize}

Prioritized predicates:

\begin{itemize}
\item Use \code{class} as the canonical entity type. Do not duplicate the same value as \code{rdf:type}.
\item Keep \code{rdf:type} only when it adds distinct KB type information not already represented by \code{class}.
\item Include user-meaningful simulator/KB predicates when present: \code{dbp:name}, \code{dbp:color}, \code{oro:isAt}, \code{oro:isOn}, \code{oro:contains}, \code{foaf:knows}.
\item Drop noisy/default prompt relations such as provenance, detector score, coordinates, backend, source, and \code{last\_seen\_*}.
\item Do not emit \code{counts}; prompts should inspect the visible \code{entities} list directly.
\item Keep raw triples available to skill-local execution paths when a skill needs exact RDF.
\end{itemize}



\subsection{Grounding Projection Pipeline}

\begin{lstlisting}[caption={Diagram source from canonical Markdown}]
flowchart LR
    A["/scene/summary"] --> P["project_llm_grounded_context"]
    B["tracked people"] --> P
    C["/kb/query rows and RDF triples"] --> P
    D["state_t0 flag (disabled by default)"] --> P
    P --> G["compact grounded_context"]
    A --> X["execution_context / skill-local raw evidence"]
    C --> X
    X --> S["AB=1 skill requery and validation"]
\end{lstlisting}

Rules:

\begin{enumerate}
\item Backend seams stay intact: \code{/scene/summary}, \code{/kb/query}, \code{/kb/revise}, tracked people, and skill action results are not rewritten by the projection.
\item The compact projection is deterministic and JSON-native.
\item Text-derived hydration is fallback-only when no structured scene/KB entities exist.
\item \code{state\_t0} remains behind \code{grounded\_context\_include\_state\_t0}; default is disabled.
\item Planner-only details such as coordinates and recency are not part of the default chatbot/planner LLM view.
\end{enumerate}



\subsection{Planner Request Envelope}

Transport-level envelope shape as serialized in the \code{Intent.data} field:

\begin{lstlisting}
{
  "request_id": "turn_123",
  "goal_id": "goal_turn_123",
  "parent_goal_id": "",
  "supersedes_goal_id": "",
  "request_kind": "new_goal",
  "goal_text": "look at the person and report completion",
  "normalized_intents": ["look_at"],
  "scene_targets": ["person"],
  "dialogue_context": ["assistant:Okay, I will do that."],
  "grounded_context": {
    "entities": [
      {
        "id": "anonymous_person_ehfbf",
        "label": null,
        "kind": "person",
        "class": "Human",
        "visible": true,
        "relations": []
      }
    ]
  },
  "planner_mode": "default",
  "dialogue_turn_id": "role:turn"
}
\end{lstlisting}

Variables:

\begin{itemize}
\item Identification and lineage: \code{request\_id}, \code{goal\_id}, \code{parent\_goal\_id}, \code{supersedes\_goal\_id}.
\item Transition control: \code{request\_kind}.
\item Planning objective: \code{goal\_text}, \code{normalized\_intents}, \code{scene\_targets}.
\item Context: \code{dialogue\_context}, \code{grounded\_context}.
\item Runtime mode: \code{planner\_mode}, \code{dialogue\_turn\_id}.
\end{itemize}

The planner prompt receives this compact payload plus execution feedback, skill registry, allowed step types, allowed skill names, and output contract.



\subsection{Planner LLM Internals}

\begin{lstlisting}[caption={Diagram source from canonical Markdown}]
flowchart TD
    R["PlannerRequest.from_payload"] --> G["PlannerGate admission already completed"]
    G --> E["PlannerEngine.plan_request"]
    E --> RB{"Rule-based simple motion?"}
    RB -- yes --> BP["Build deterministic rule plan"]
    RB -- no --> PR["Build LLM prompt payload"]
    PR --> L["LLM provider"]
    L --> V["Parse JSON + validate steps"]
    V -->|valid| OUT["build_plan_payload"]
    V -->|invalid| RETRY["validation retry prompt"]
    RETRY --> V
    V -->|still invalid| FAIL["fail/clarify payload"]
    BP --> OUT
    OUT --> I["/intents"]
\end{lstlisting}

Planner prompt payload contains:

\begin{itemize}
\item \code{request}: compact planner request.
\item \code{goal\_id}, \code{plan\_version}.
\item optional \code{state\_t0} only when explicitly enabled by flag.
\item \code{execution\_feedback} during replans.
\item \code{skill\_registry}, \code{allowed\_step\_types}, \code{allowed\_skill\_names}, \code{allowed\_motion\_objects}.
\item \code{output\_contract}.
\end{itemize}

Planner prompt payload does not contain:

\begin{itemize}
\item raw \code{user\_text} in normal operation.
\item raw detector coordinates or scores.
\item full RDF dumps.
\item removed request seams.
\end{itemize}



\subsection{Planner Output}

\begin{lstlisting}
{
  "plan": {
    "goal_id": "goal_turn_123",
    "plan_id": "plan_123",
    "plan_version": 2,
    "status": "planning",
    "validation_status": "draft",
    "user_facing_reason": "",
    "replan_hint": "",
    "retry_budget": 1,
    "scene_targets": ["person"],
    "context_ref": {
      "captured_at_sec": 0.0,
      "observer": "",
      "backend": ""
    },
    "communication_policy": {
      "emit_acknowledge": false,
      "emit_progress": false,
      "emit_completion": true,
      "emit_failure": true
    },
    "communication_policy_source": "planner_engine:plan",
    "steps": [
      {
        "id": "step_1",
        "type": "skill",
        "name": "scan",
        "args": {},
        "requires": [],
        "on_failure": "replan",
        "retry_budget": 0
      },
      {
        "id": "step_2",
        "type": "skill",
        "name": "report_result",
        "args": {},
        "requires": ["step_1"],
        "on_failure": "fail",
        "retry_budget": 0
      }
    ]
  }
}
\end{lstlisting}

Variables:

\begin{itemize}
\item Plan lineage: \code{goal\_id}, \code{plan\_id}, \code{plan\_version}.
\item Plan state: \code{status}, \code{validation\_status}, \code{user\_facing\_reason}, \code{replan\_hint}.
\item Failure detail: \code{failure\_reason} is present only when there is an actual failure, invalid output, or terminal blocked condition.
\item Resilience: \code{retry\_budget}.
\item Traceability: \code{scene\_targets}, \code{context\_ref}.
\item Communication: \code{communication\_policy}, \code{communication\_policy\_source}.
\item Execution content: \code{steps[]} with \code{id}, \code{type}, \code{name}, \code{args}, \code{requires}, \code{on\_failure}, \code{retry\_budget}.
\end{itemize}



\subsection{Communication Policy}

\code{communication\_policy} is planner dialogue-act permission, not direct speech execution.

\begin{itemize}
\item \code{emit\_acknowledge}: allow a planner-owned acknowledgement dialogue act on plan acceptance.
\item \code{emit\_progress}: allow planner-owned progress dialogue acts on step start/progress.
\item \code{emit\_completion}: allow planner-owned completion dialogue acts on plan completion.
\item \code{emit\_failure}: allow planner-owned failure/help dialogue acts on execution failures.
\item \code{communication\_policy\_source}: deterministic provenance showing where the final policy was resolved, usually \code{planner\_engine:<mode>}.
\end{itemize}

Defaults come from \code{planner\_common.normalize\_communication\_policy}. The planner may override policy flags in its JSON output; the shared contract normalizes them before supervisor use.



\subsection{Supervisor Lifecycle}

\begin{lstlisting}[caption={Diagram source from canonical Markdown}]
stateDiagram-v2
    [*] --> idle
    idle --> planning: accepted request
    planning --> executing: valid plan
    planning --> waiting_user: clarify
    planning --> failed: invalid/fail
    executing --> replanning: step failed + retry budget
    executing --> waiting_user: needs user input
    executing --> completed: all steps succeeded
    executing --> failed: terminal failure
    replanning --> executing: new valid plan
    replanning --> waiting_user: no safe continuation
    waiting_user --> planning: clarification_answer
    completed --> idle
    failed --> idle
\end{lstlisting}

Supervisor responsibilities:

\begin{enumerate}
\item Track \code{goal\_id}, \code{plan\_id}, and \code{plan\_version}.
\item Admit new plans and suppress duplicate plan versions.
\item Dispatch only executable steps through the orchestrator.
\item Consume execution feedback and decrement retry budget.
\item Ask for replans when failure policy and budget allow.
\item Emit planner dialogue acts for clarification, progress, completion, or failure according to \code{communication\_policy}.
\end{enumerate}



\subsection{Scan and Report Evidence Flow}

\begin{lstlisting}[caption={Diagram source from canonical Markdown}]
flowchart TD
    A["User: what do you see?"] --> B["chatbot_llm compact grounded_context"]
    B --> C["planner_llm plans scan + report_result"]
    C --> D["nao_orchestrator dispatches scan"]
    D --> E["scan skill requeries scene/KB"]
    E --> F["scan result payload with summary_text"]
    F --> G["orchestrator stores latest result"]
    G --> H["report_result step with empty args"]
    H --> I["orchestrator reuses latest live result"]
    I --> J["execution_feedback to planner"]
\end{lstlisting}

Critical guardrail:

\begin{itemize}
\item The planner must not prefill \code{report\_result.args.summary\_text} after \code{scan}.
\item The scan step owns fresh perception.
\item \code{report\_result} with empty args tells the orchestrator to report the latest live skill result.
\item \code{summary\_text} is only allowed for known non-perception facts or already completed non-perception results.
\end{itemize}

This prevents a stale prompt snapshot from being mistaken for a fresh scan.



\subsection{Execution Feedback}

\begin{lstlisting}
{
  "goal_id": "goal_turn_123",
  "plan_id": "plan_123",
  "plan_version": 2,
  "intent": "raw_user_input",
  "source": "planner_llm",
  "event_type": "step_succeeded",
  "status": "succeeded",
  "reason": "",
  "validation_status": "valid",
  "replan_hint": "",
  "retry_budget": 1,
  "blocking": false,
  "unmet_preconditions": [],
  "needs_user_input": false,
  "scene_targets": ["person"],
  "validation_errors": [],
  "timestamp_sec": 1777040000.0,
  "result_summary": "Look-at completed for person_1",
  "result_payload": {
    "skill": "look_at",
    "target": "person_1",
    "target_kind": "people",
    "target_found": true,
    "summary_text": "Look-at completed for person_1",
    "confidence_policy": "grounded_current_observation"
  },
  "step": {
    "id": "step_1",
    "type": "skill",
    "name": "look_at",
    "retry_budget": 0,
    "on_failure": "replan",
    "requires": []
  }
}
\end{lstlisting}

Evidence fields:

\begin{itemize}
\item \code{result\_summary}: compact human-readable evidence summary.
\item \code{result\_payload}: structured skill result.
\item \code{scene\_targets}: targets used for replan context.
\item \code{blocking}, \code{needs\_user\_input}, \code{unmet\_preconditions}: supervisor decision hints.
\end{itemize}



\subsection{Planner Dialogue Act}

\begin{lstlisting}
{
  "goal_id": "goal_turn_123",
  "plan_id": "plan_123",
  "plan_version": 2,
  "act": "ask_clarification",
  "priority": "normal",
  "await_user_response": true,
  "reason": "missing target object",
  "text_hint": "Which person should I look at?",
  "slots_needed": ["target"],
  "context": {
    "scene_targets": ["person"],
    "status": "waiting_user"
  }
}
\end{lstlisting}

Allowed \code{act} values:

\begin{itemize}
\item \code{acknowledge}
\item \code{progress\_update}
\item \code{ask\_clarification}
\item \code{ask\_for\_help}
\item \code{explain\_failure}
\item \code{notify\_completion}
\item \code{notify\_cancellation}
\end{itemize}

Dialogue acts travel through \code{nao\_orchestrator} to \code{dialogue\_manager}. The planner does not speak directly.



\subsection{Chatbot / Planner / Skill Ownership Map}

\begin{lstlisting}[caption={Diagram source from canonical Markdown}]
flowchart LR
    C["chatbot_llm\nuser-facing grounding + planner handoff"] --> P["planner_llm\nplan + supervise"]
    P --> O["nao_orchestrator\nadmit + dispatch + feedback"]
    O --> S1["scan\nfresh perception"]
    O --> S2["look_at\ntarget resolution"]
    O --> S3["navigate_to\nmovement evidence"]
    O --> S4["report_result\nlatest result speech payload"]
    S1 --> K["scene/KB/tracked people"]
    S2 --> K
    S3 --> K
    S4 --> O
    O --> D["dialogue_manager\nspeech realization"]
\end{lstlisting}

Methodology:

\begin{enumerate}
\item Keep LLM context compact and semantically rich.
\item Keep raw sensor/KB evidence out of every prompt by default.
\item Validate planner output outside the LLM.
\item Require skills to requery live state before reporting execution success.
\item Use planner dialogue acts for coordination, not direct speech generation.
\end{enumerate}



\subsection{Replanning and Join Semantics}

\begin{lstlisting}[caption={Diagram source from canonical Markdown}]
flowchart TD
    A["Receive /planner/request"] --> B{"Existing active goal?"}
    B -- No --> C["Create new goal lineage"]
    B -- Yes --> D{"request_kind"}
    D -- new_goal --> E["Supersede or reject based on lineage"]
    D -- goal_update --> F["Keep goal_id continuity"]
    D -- clarification_answer --> F
    D -- cancel_request --> G["Cancel active goal"]
    C --> H["Publish /intents"]
    E --> H
    F --> H
    H --> I["Execute steps"]
    I --> J["Publish execution feedback"]
    J --> K{"Need replan?"}
    K -- Yes --> L["Increment plan_version and replan"]
    K -- No --> M["Continue or complete"]
    L --> H
\end{lstlisting}

Current policy summary:

\begin{enumerate}
\item Single active-goal queue model.
\item \code{goal\_id} keeps goal continuity.
\item \code{plan\_id} + \code{plan\_version} encode plan lineage.
\item Mid-join when active \code{step\_id} is present in the new plan.
\item Front-join from first pending step when mid-join cannot be resolved.
\item Duplicate suppression compares \code{goal\_id}, \code{plan\_id}, and \code{plan\_version}.
\end{enumerate}



\subsection{Live Supervisor Demo Script}

\begin{enumerate}
\item Trigger one command and show chatbot output (\code{route=execution}).
\item Show compact \code{grounded\_context} on the planner request.
\item Show admitted request on \code{/planner/request}.
\item Show planner output on \code{/intents} and highlight lineage fields.
\item Run one scan/report case and show that \code{report\_result.args} is empty.
\item Run one failure/clarification path.
\item Show execution feedback payload and resulting dialogue act payload.
\item Close with ownership map: chatbot grounds, planner plans, orchestrator dispatches, skills provide live evidence, dialogue manager speaks.
\end{enumerate}
\end{document}